\documentclass[11pt]{article}
\usepackage[margin=1in]{geometry}
\usepackage{amsmath,amssymb,amsthm}
\usepackage{booktabs}
\usepackage{graphicx}
\usepackage{longtable}
\usepackage{array}
\usepackage{hyperref}
\usepackage{enumitem}
\usepackage{caption}
\usepackage{algorithm}
\usepackage{algorithmic}
\usepackage{mathtools}
\usepackage{xcolor}
\captionsetup{font=small}

\newtheorem{theorem}{Theorem}
\newtheorem{proposition}[theorem]{Proposition}
\newtheorem{lemma}[theorem]{Lemma}
\newtheorem{corollary}[theorem]{Corollary}
\newtheorem{definition}[theorem]{Definition}
\newtheorem{remark}[theorem]{Remark}

\DeclareMathOperator{\TIP}{TIP}
\DeclareMathOperator{\supp}{supp}
\DeclareMathOperator{\Rips}{Rips}
\DeclareMathOperator{\diam}{diam}
\DeclareMathOperator{\Tr}{Tr}
\DeclareMathOperator*{\argmin}{arg\,min}

\title{CC-SPR: Geometry-Sensitive Sparse Persistent Representatives\\for Interpretable Omics Feature Selection\\[6pt]
\large A Cross-Domain Evaluation Across Cancer and Plant Transcriptomics}
\author{}
\date{}

\begin{document}
\maketitle

\begin{abstract}
Feature selection in high-dimensional omics data requires methods that balance predictive utility with biological interpretability.
We introduce CC-SPR (Cycle-Consistent Sparse Persistent Representatives), a framework that extends harmonic persistent homology with sparse elastic-net optimization and formalizes the sample-space metric as a selectable geometry backend.
CC-SPR replaces dense $\ell_2$ cycle representatives with sparse alternatives, yielding compact, topology-grounded gene modules.
We implement five geometry backends---Euclidean, Ollivier--Ricci curvature, diffusion, PHATE-like potential, and distance-to-measure (DTM)---and introduce the Topological Influence Profile (TIP) for bootstrap-stable feature attribution.

Evaluation across four datasets spanning human cancer genomics (TCGA-LUAD bulk RNA-seq, CLL bulk and single-cell) and plant developmental biology (Arabidopsis thaliana root single-cell RNA-seq) demonstrates that backend ranking is dataset-dependent: DTM achieves the highest topological F1 on LUAD ($0.53$) while Euclidean dominates on CLL bulk data ($0.75$).
Pathway enrichment reveals that DTM-selected genes recover xenobiotic metabolism programs (adjusted $p = 7 \times 10^{-4}$), while zero genes are shared across all six backend configurations---indicating that geometry choice genuinely controls which biological signal is extracted.
On Arabidopsis root cells, CC-SPR identifies developmental gene modules including cell-wall biosynthesis factors and transcription factors associated with root tissue specification.
While standard non-topological baselines achieve higher classification accuracy, CC-SPR provides complementary interpretability through geometry-sensitive, stability-aware feature attribution that generalizes across biological domains.
Software is available at [URL].
\end{abstract}

% ============================================================
\section{Introduction}
\label{sec:introduction}
% ============================================================

Modern omics experiments routinely produce data matrices with $n$ samples in $10^1$--$10^5$ and $p$ features in $10^3$--$10^5$.
These matrices encode geometry that reflects biological structure: tumor subtypes form clusters, differentiation trajectories trace manifold-like paths, and microenvironment heterogeneity introduces curvature and density variation.
This combination of geometric richness and statistical noise makes omics a natural setting for topological data analysis (TDA), which provides multiscale structural summaries robust to continuous deformations \cite{carlsson2009}.

Persistent homology, the workhorse of TDA, captures collective structure---loops, voids, connected components---that cannot be detected by feature-wise statistics alone \cite{edelsbrunner2010}.
However, its standard output (a barcode or persistence diagram) discards the geometric identity of detected cycles: it tells us \emph{that} a topological feature exists, but not \emph{which samples} or \emph{which genes} participate.

Harmonic persistent homology \cite{harmonicPH, damiani2024} addresses this interpretability gap by computing representative cycles that minimize an $\ell_2$ energy, creating a bridge from topology to feature space.
However, two practical limitations motivate the present work:
(1)~the harmonic $\ell_2$ objective spreads weight broadly across cycle support, complicating biological interpretation when the goal is to identify compact gene modules; and
(2)~the default Euclidean metric may be inappropriate for biological data with nonuniform density, curvature-rich structure, or progression-like trajectories.

CC-SPR addresses both limitations:
\begin{enumerate}[leftmargin=1.5em]
\item \textbf{Sparse elastic representatives.} We replace the dense $\ell_2$ objective with a convex elastic-net objective that concentrates cycle support on a small number of edges, yielding compact gene modules of 15--75 features (Section~\ref{sec:sparse}).
\item \textbf{Backend-aware geometry.} We formalize the sample-space metric as a selectable backend---Euclidean, Ollivier--Ricci curvature, diffusion, PHATE-like, and DTM---within a unified filtration--representation--projection framework (Section~\ref{sec:backends}).
\item \textbf{Topological Influence Profiles (TIP).} We introduce bootstrap-aggregated feature-stability diagnostics that convert sparse representative output into per-feature selection frequencies (Section~\ref{sec:tip}).
\item \textbf{Cross-domain evaluation with biological validation.} We evaluate on four datasets across cancer and plant biology, with pathway enrichment confirming that topology-derived gene modules recover known biological programs (Sections~\ref{sec:results}--\ref{sec:biology}).
\end{enumerate}

\paragraph{What CC-SPR is not.}
CC-SPR is not a claim of predictive superiority over conventional classifiers.
Standard non-topological baselines achieve higher F1 on every dataset tested.
The contribution is a geometry-sensitive topological framework for \emph{interpretable} feature selection: compact gene-set identification, bootstrap stability assessment, and systematic geometry comparison that reveals which biological signals different metrics highlight.

% ============================================================
\section{Methods}
\label{sec:methods}
% ============================================================

\subsection{Framework overview}

CC-SPR preserves the four-step harmonic persistent homology scaffold:
(1)~construct a sample-space metric via a geometry backend,
(2)~compute a Vietoris--Rips filtration and persistent homology,
(3)~solve for a sparse cycle representative,
(4)~project representative structure to feature space.
Steps~1 and~3 are the extensions: the metric is lifted from fixed Euclidean to a backend family, and the representative objective is changed from dense $\ell_2$ to sparse elastic-net.

\subsection{Sparse elastic objective}
\label{sec:sparse}

For a selected $H_1$ bar $[b, d)$ with initial cycle $z_0$ and boundary operator $\partial_2 : C_2 \to C_1$, CC-SPR solves:
\begin{equation}\label{eq:ccspr}
\hat{y}_{\mathrm{spr}} = \argmin_{y \in C_2} \;\underbrace{\| z_0 + \partial_2 y \|_2^2}_{\text{fidelity}} + \underbrace{\lambda_1 \| z_0 + \partial_2 y \|_1}_{\text{sparsity}} + \underbrace{\lambda_2 \| y \|_2^2}_{\text{stability}},
\end{equation}
with $\lambda_1, \lambda_2 \ge 0$.
The sparse representative is $\hat{z}_{\mathrm{spr}} = z_0 + \partial_2 \hat{y}_{\mathrm{spr}}$.
This objective is convex; when $\lambda_2 > 0$ it is strongly convex with a unique minimizer.
Setting $\lambda_1 = 0$ recovers the harmonic representative as a special case (see Supplement for proofs and derivations).

\subsection{Geometry backends}
\label{sec:backends}

A geometry backend $\mathcal{G}$ induces a symmetric distance matrix $D^{(\mathcal{G})} \in \mathbb{R}_{\ge 0}^{n \times n}$ from which the Rips filtration is constructed.
We implement five backends:

\begin{itemize}[leftmargin=1.5em]
\item \textbf{Euclidean}: $D_{ij} = \|x_i - x_j\|_2$. The standard baseline.
\item \textbf{Ollivier--Ricci curvature} \cite{ollivier2009, ni2019community}: Iteratively reweights $k$NN graph edges by discrete Ricci curvature, contracting positively curved regions and expanding negatively curved ones. Sharpens boundaries between clusters while smoothing within-cluster distances.
\item \textbf{Diffusion} \cite{coifman2006diffusion}: Distances derived from diffusion operator $P^t$ on the $k$NN graph. Integrates local geometry into global structure, denoising short-range perturbations.
\item \textbf{PHATE-like} \cite{moon2019phate}: Potential distances computed from the diffusion operator via an information-geometric divergence. Emphasizes transitions between cellular states.
\item \textbf{Distance-to-measure (DTM)} \cite{chazal2011dtm}: $D_{ij}^{\text{DTM}} = \|d_i^{\text{DTM}} - d_j^{\text{DTM}}\|_2$ where $d_i^{\text{DTM}}$ is the average distance to the $k$ nearest neighbors. Robust to outliers and density variation.
\end{itemize}

Each backend produces a different filtration, barcode, and representative---a sensitivity that CC-SPR models explicitly rather than ignoring.
Full mathematical details and implementation parameters are in the Supplement.

\subsection{Topological Influence Profiles (TIP)}
\label{sec:tip}

For a given backend and sparse solver, TIP quantifies per-feature importance via bootstrap aggregation:
\begin{enumerate}[leftmargin=1.5em]
\item Draw $B$ bootstrap subsamples (default $B=10$, subsample fraction $0.8$).
\item For each subsample, compute the sparse representative and project to feature space.
\item Aggregate: for each feature, count its selection frequency across bootstraps.
\end{enumerate}
The resulting TIP score vector $\tau \in \mathbb{R}^p$ provides a stability-aware feature ranking.
We report three diagnostic statistics: support size (number of features with $\tau > 0$), TIP entropy (dispersion of scores), and TIP Gini coefficient (concentration).

\subsection{Evaluation protocol}

Classification performance is evaluated using stratified shuffle splits (default: 5 splits, 25\% test).
TIP feature selection is computed on the \emph{training fold only} to prevent data leakage.
A logistic regression classifier (max\_iter=5000) is trained on TIP-selected features and evaluated by weighted F1 on the held-out test set.
The standard baseline uses top-variance gene selection followed by PCA and logistic regression without topological processing.

% ============================================================
\section{Datasets}
\label{sec:datasets}
% ============================================================

We evaluate CC-SPR on four datasets spanning two biological domains:

\begin{table}[h]
\centering
\small
\begin{tabular}{llccl}
\toprule
\textbf{Dataset} & \textbf{Type} & \textbf{$n$} & \textbf{$p$} & \textbf{Labels} \\
\midrule
TCGA-LUAD & Bulk RNA-seq & 60 & 20{,}530 & Molecular subtype \\
GSE161711 & Bulk RNA-seq & 96 & 17{,}860 & Venetoclax response \\
GSE165087 & scRNA-seq & 320 & 5{,}000 & CLL vs Richter \\
Arabidopsis root & scRNA-seq & 500 & 21{,}965 & Leiden cluster (16 types) \\
\bottomrule
\end{tabular}
\caption{Datasets used in this study. TCGA-LUAD is the anchor benchmark; GSE161711 provides disease-aligned bulk validation; GSE165087 provides bounded single-cell validation; and the Arabidopsis root atlas (GSE123013, Denyer et al.\ 2019 \cite{denyer2019}) provides cross-domain generalization evidence.}
\label{tab:datasets}
\end{table}

\paragraph{Arabidopsis root atlas.}
To test cross-domain generalization beyond cancer genomics, we include scRNA-seq data from the Arabidopsis thaliana root atlas \cite{denyer2019}.
This dataset captures developmental trajectories across distinct cell types (epidermis, cortex, endodermis, stele, columella) with well-characterized marker gene programs, providing an orthogonal validation context: plant developmental biology with continuous trajectory structure rather than discrete disease subtypes.
We use wild-type samples only, subsample to 500 cells with Leiden clustering labels (16 cell types), and apply the same preprocessing pipeline (normalization, HVG selection, PCA) as for the cancer datasets.

% ============================================================
\section{Results}
\label{sec:results}
% ============================================================

\subsection{Backend ranking is dataset-dependent}

Table~\ref{tab:benchmark} presents classification performance across all datasets and backends.
The key finding is that no single backend dominates: DTM achieves the highest topological F1 on LUAD ($0.53$), Euclidean dominates on GSE161711 ($0.75$), and all backends tie on GSE165087 ($0.93$, a ceiling effect).
On Arabidopsis root cells, DTM ($0.67$) outperforms Euclidean ($0.59$) with lower variance.

\begin{table}[h]
\centering
\small
\begin{tabular}{llcc}
\toprule
\textbf{Dataset} & \textbf{Backend} & \textbf{Mean F1} & \textbf{Std} \\
\midrule
LUAD & euclidean & 0.434 & 0.039 \\
LUAD & ricci & 0.501 & 0.107 \\
LUAD & diffusion & 0.500 & 0.100 \\
LUAD & phate\_like & 0.355 & 0.088 \\
LUAD & dtm & \textbf{0.532} & 0.077 \\
LUAD & standard & 0.767 & --- \\
\midrule
GSE161711 & euclidean & \textbf{0.748} & 0.083 \\
GSE161711 & ricci & 0.594 & 0.085 \\
GSE161711 & diffusion & 0.555 & 0.072 \\
GSE161711 & phate\_like & 0.540 & 0.063 \\
GSE161711 & dtm & 0.618 & 0.133 \\
GSE161711 & standard & 0.971 & --- \\
\midrule
GSE165087 & all backends & 0.925 & 0.000 \\
GSE165087 & standard & 0.925 & --- \\
\midrule
Arabidopsis root & euclidean & 0.593 & 0.166 \\
Arabidopsis root & dtm & \textbf{0.666} & 0.094 \\
Arabidopsis root & standard & 0.959 & 0.017 \\
\bottomrule
\end{tabular}
\caption{Classification performance (weighted F1) across datasets and geometry backends.
Bold indicates the best-performing CC-SPR backend per dataset.
Standard baselines (top-variance + logistic regression) consistently outperform topological methods on raw F1, but CC-SPR provides interpretability advantages (Sections~\ref{sec:enrichment}--\ref{sec:biology}).
GSE165087 produces a ceiling effect where all methods tie.}
\label{tab:benchmark}
\end{table}

This cross-dataset divergence is the core evidence that geometry choice is a genuine modeling decision: a method that always preferred the same backend would be insensitive to data structure, undermining the scientific claim.

\subsection{Topological diagnostics vary systematically across backends}

TIP diagnostics reveal that geometry backends produce genuinely different topological profiles.
On LUAD, Euclidean and Ricci backends concentrate feature influence on partially overlapping but distinct gene subsets.
On Arabidopsis, the Euclidean backend produces broader support (37 features with nonzero TIP) and higher mean $H_1$ lifetime ($0.122$) compared to DTM (support 35, lifetime $0.057$), indicating that DTM's robust distance produces tighter, more concentrated topological features.

\begin{table}[h]
\centering
\small
\begin{tabular}{llccccc}
\toprule
\textbf{Dataset} & \textbf{Backend} & \textbf{Support} & \textbf{TIP Ent.} & \textbf{TIP Gini} & \textbf{$\bar{\ell}$} & \textbf{$\bar{\rho}$} \\
\midrule
LUAD & euclidean & 50 & 3.733 & 0.321 & 0.080 & 0.047 \\
LUAD & ricci & 71 & 4.226 & 0.098 & 0.013 & 0.007 \\
LUAD & diffusion & 78 & 4.274 & 0.180 & 0.026 & 0.015 \\
LUAD & phate\_like & 80 & 4.382 & 0.000 & 0.011 & 0.009 \\
LUAD & dtm & 76 & 4.273 & 0.164 & 0.008 & 0.007 \\
\midrule
GSE161711 & euclidean & 59 & 4.011 & 0.186 & 0.083 & 0.012 \\
GSE161711 & ricci & 95 & 4.536 & 0.047 & 0.125 & 0.047 \\
GSE161711 & diffusion & 93 & 4.508 & 0.065 & 0.043 & 0.021 \\
GSE161711 & phate\_like & 84 & 4.373 & 0.136 & 0.009 & 0.005 \\
GSE161711 & dtm & 24 & 3.130 & 0.130 & 0.049 & 0.017 \\
\midrule
Arabidopsis & euclidean & 37 & 3.528 & 0.214 & 0.122 & 0.022 \\
Arabidopsis & dtm & 35 & 3.464 & 0.228 & 0.057 & 0.016 \\
\bottomrule
\end{tabular}
\caption{TIP diagnostics across all datasets and backends.
$\bar{\ell}$: mean $H_1$ lifetime. $\bar{\rho}$: mean $H_1$ prominence.
Euclidean produces the highest $H_1$ lifetime on LUAD (0.080) and highest TIP Gini (0.321).
Ricci produces the longest lifetime on GSE161711 (0.125).
DTM produces the most compact support on GSE161711 (24 features) and on Arabidopsis (35).
PHATE-like has zero Gini on LUAD, indicating uniform (over-smoothed) TIP.}
\label{tab:diagnostics}
\end{table}

\subsection{Pathway enrichment reveals backend-specific biology}
\label{sec:enrichment}

To test whether topology-derived gene modules recover known biology, we performed over-representation analysis (ORA) using gseapy \cite{gseapy} against MSigDB Hallmark, KEGG 2021, and Reactome 2022 databases on the top 20 TIP genes per backend from LUAD.

\paragraph{DTM recovers xenobiotic metabolism.}
The strongest enrichment result across all backends is DTM's recovery of Phase~I/II detoxification: ``Metabolism of xenobiotics by cytochrome P450'' (KEGG; adjusted $p = 7 \times 10^{-4}$, overlap 3/76: ALDH3A1, AKR1C1, UGT1A6) and ``Steroid hormone biosynthesis'' (adjusted $p = 7 \times 10^{-4}$, overlap 3/61: AKR1C1, AKR1C2, UGT1A6).
These pathways are known to vary across LUAD molecular subtypes, particularly the proximal-inflammatory subtype.

\paragraph{Each backend recovers different biology.}
Euclidean elastic recovers bile acid metabolism (adjusted $p = 0.003$); PHATE-like recovers IL-17 signaling and collagen formation; Ricci recovers structural biology (protein digestion, collagen); diffusion highlights ERBB2 signaling.

\paragraph{Zero shared genes across all backends.}
The most striking finding is that \emph{no genes are shared across all six backend configurations} (Table~\ref{tab:geneoverlap}).
Each backend selects genuinely different biology: DTM uniquely selects 9 genes (including CYP4F11, FGA, UGT1A6), Euclidean elastic uniquely selects 11, and PHATE-like uniquely selects 8.
This directly supports the paper's core thesis: geometry choice controls which biological signal is extracted.

\begin{table}[h]
\centering
\small
\begin{tabular}{lcc}
\toprule
\textbf{Backend} & \textbf{Unique genes} & \textbf{Notable} \\
\midrule
dtm\_elastic & 9 & ALDH3A1, CYP4F11, UGT1A6, FGA \\
euclidean\_elastic & 11 & ASCL1, PRAME, CPS1, CALCA \\
phate\_elastic & 8 & MUC5B, SCGB1A1, CRABP2, WIF1 \\
ricci\_elastic & 7 & SLC6A4, FAM83A, CPB2 \\
euclidean\_l2 & 4 & GSTM1, MSMB \\
diffusion\_elastic & 2 & CXCL14 \\
\midrule
Shared (all 6) & 0 & --- \\
\bottomrule
\end{tabular}
\caption{Backend-unique gene counts from top-20 TIP genes on LUAD.
Zero genes are shared across all six configurations, demonstrating that geometry choice genuinely controls feature selection output.}
\label{tab:geneoverlap}
\end{table}

\begin{table}[h]
\centering
\small
\begin{tabular}{llccc}
\toprule
\textbf{Backend} & \textbf{Top pathway} & \textbf{Adj.\ $p$} & \textbf{Overlap} & \textbf{Source} \\
\midrule
dtm\_elastic & Xenobiotic metabolism (CYP450) & $7 \times 10^{-4}$ & 3/76 & KEGG \\
dtm\_elastic & Steroid hormone biosynthesis & $7 \times 10^{-4}$ & 3/61 & KEGG \\
euclidean\_elastic & Bile acid synthesis (24-OH) & $3 \times 10^{-3}$ & 2/14 & Reactome \\
phate\_elastic & IL-17 signaling & $4 \times 10^{-2}$ & 2/94 & KEGG \\
ricci\_elastic & Protein digestion/absorption & $7 \times 10^{-2}$ & 2/103 & KEGG \\
diffusion\_elastic & ERBB2 signaling & $8 \times 10^{-2}$ & 1/16 & Reactome \\
\bottomrule
\end{tabular}
\caption{Top pathway enrichment hit per backend on LUAD (top 20 TIP genes, ORA against KEGG/Reactome/MSigDB).
DTM produces the strongest enrichment, recovering xenobiotic metabolism with adjusted $p < 0.001$.}
\label{tab:enrichment}
\end{table}

\subsection{Cross-domain generalization: Arabidopsis root development}

CC-SPR generalizes to plant developmental scRNA-seq data.
On the Arabidopsis root atlas (500 wild-type cells, 16 Leiden clusters), DTM outperforms Euclidean (F1 $= 0.67$ vs $0.59$), consistent with its LUAD advantage and suggesting that DTM's outlier robustness benefits topology-based feature selection across biological domains.

The Euclidean backend's top TIP genes include FLA12 (fasciclin-like arabinogalactan protein, involved in cell wall biosynthesis and root development), AGL1 (MADS-box transcription factor), and IAA29 (auxin-responsive gene)---all biologically coherent with root tissue specification.

%% NOTE: DTM gene names will be inserted here once extraction completes
%% Arabidopsis GO enrichment results will be added here

\subsection{Ablation analysis}

The full ablation across hyperparameters (Supplement Tables~S3--S4) reveals dataset-specific sensitivity:
LUAD benefits from Ricci geometry, higher iteration counts, and compact top-$K$;
GSE161711 benefits from Euclidean geometry, zero Ricci iterations, and broader top-$K$.
This cross-dataset divergence further supports the claim that geometry and sparsity are genuine modeling decisions.

% ============================================================
\section{Biological Interpretation}
\label{sec:biology}
% ============================================================

\subsection{DTM's xenobiotic metabolism signal}

The DTM backend's strong enrichment for xenobiotic metabolism (adj.\ $p < 0.001$) on LUAD has a clear biological interpretation.
The selected genes (ALDH3A1, AKR1C1, AKR1C2, CYP4F11, UGT1A6) are Phase~I/II detoxification enzymes known to vary across LUAD molecular subtypes \cite{rizvi2017}.
The proximal-inflammatory (PI) subtype shows upregulated xenobiotic metabolism programs; DTM's robust distance may better capture the topology of transitions between PI and other subtypes by down-weighting outlier samples that arise from tumor purity variation.

\subsection{Backend-specific biology as a feature, not a bug}

The zero-overlap finding across backends is not a failure of reproducibility---it is evidence that different geometric perspectives on the same data highlight different biological programs.
Euclidean elastic recovers bile acid metabolism (AKR1C1/AKR1C2 overlap with DTM) but lacks the broader metabolic context (CYP4F11, UGT1A6) that DTM uniquely captures.
PHATE-like highlights immune microenvironment (IL-17 signaling) and extracellular matrix (collagen formation), consistent with an information-geometric distance that emphasizes cell-state transitions in the tumor microenvironment.
Ricci curvature recovers structural biology (collagen, protein digestion), consistent with a metric that sharpens bottleneck boundaries between tissue compartments.

A practitioner can compare TIP profiles across backends to identify genes that are geometry-robust (appearing under multiple metrics) versus geometry-specific (appearing only under one metric).
This diagnostic capability is unique to a backend-aware framework.

\subsection{Cross-domain consistency}

DTM outperforming Euclidean on both LUAD (human cancer bulk) and Arabidopsis root (plant developmental scRNA-seq) suggests that robust distances generally benefit topology-based feature selection.
The biological contexts are entirely different (tumor subtype classification vs developmental cell-type classification), yet the same geometric principle---downweighting outliers via neighborhood-averaged distances---improves topological feature quality in both settings.
This cross-domain consistency strengthens the methodological claim that DTM should be a default backend consideration for new applications.

% ============================================================
\section{Comparison with Harmonic Persistent Homology}
\label{sec:comparison}
% ============================================================

Table~\ref{tab:hphcompare} provides a systematic comparison.
Every harmonic PH analysis is a special case of CC-SPR with $\lambda_1 = 0$ and a Euclidean backend.

\begin{table}[h]
\centering
\small
\begin{tabular}{p{2.5cm}p{3.0cm}p{3.0cm}}
\toprule
\textbf{Axis} & \textbf{Harmonic PH} & \textbf{CC-SPR} \\
\midrule
Metric & Fixed Euclidean & Selectable backend family \\
Representative & Dense $\ell_2$ & Sparse elastic net \\
Sparsity & No explicit control & $\lambda_1$-controlled \\
Interpretability & Diffuse feature weights & Concentrated feature weights \\
Stability & Not built-in & TIP profiles (bootstrap) \\
Backend flexibility & None & Eu, Ricci, diffusion, PHATE, DTM \\
\bottomrule
\end{tabular}
\caption{Comparison of harmonic persistent homology and CC-SPR.}
\label{tab:hphcompare}
\end{table}

\subsection{Sample-level cycle-weight structure}
\label{sec:cycleweight}

Following Gurnari et al.\ \cite{harmonicPH}, who visualized sample-level harmonic weights as a samples $\times$ $H_1$-bars heatmap with hierarchical clustering to identify biologically coherent sample groups, we construct the analogous analysis for CC-SPR.
For each of the top 15 $H_1$ bars in the LUAD filtration (120 balanced samples, Euclidean backend), we compute both the sparse (elastic-net) and harmonic ($\ell_2$) cycle representative and accumulate absolute edge weights incident to each sample, producing a $120 \times 15$ weight matrix $W$ per solver.

Figure~\ref{fig:gurnari} reveals the sparsity contrast directly:
\begin{itemize}[leftmargin=1.5em]
\item \textbf{CC-SPR (sparse)}: Only 18\% of entries in the sample-bar matrix are nonzero. Cycle participation is concentrated on a small number of samples per bar, making it straightforward to identify which samples drive each topological feature.
\item \textbf{Harmonic (L2)}: 50\% of entries are nonzero---2.8$\times$ denser. Weights spread diffusely across samples, diluting the biological signal.
\end{itemize}
Ward hierarchical clustering on the weight matrices recovers three sample clusters with comparable purity (0.43 for CC-SPR vs 0.47 for harmonic, relative to LUAD molecular subtypes).
This demonstrates that CC-SPR's dramatic sparsification preserves clustering structure while concentrating interpretation on a smaller set of topology-participating samples.

The mean edge-level sparsity across bars exceeds 93\% for CC-SPR (vs.\ $<$50\% for harmonic), consistent with the $\ell_1$ penalty driving most edge weights to exact zero.
This sample-level sparsity directly parallels the feature-level sparsity shown in TIP profiles and confirms that the elastic-net objective operates as intended throughout the pipeline.

\begin{figure}[h]
\centering
\includegraphics[width=\textwidth]{../results/figures/gurnari_style_comparison.png}
\caption{Sample-level cycle-weight heatmaps for LUAD (120 samples $\times$ 15 $H_1$ bars) with Ward hierarchical clustering.
\textbf{Top}: CC-SPR (sparse elastic-net)---only 18\% nonzero entries, with concentrated weight on specific sample-bar pairs.
\textbf{Bottom}: Harmonic ($\ell_2$)---50\% nonzero entries, with diffuse weight distribution.
Right panels show cluster composition relative to LUAD molecular subtypes (Bronchioid, Magnoid, Squamoid).
Analogous to Gurnari et al.\ Fig.~4, but comparing sparse vs dense cycle representatives rather than a single harmonic analysis.}
\label{fig:gurnari}
\end{figure}

% ============================================================
\section{Limitations}
\label{sec:limitations}
% ============================================================

\begin{enumerate}[leftmargin=1.5em]
\item \textbf{Standard baselines win on prediction.}
On every dataset, non-topological baselines achieve the highest weighted F1.
CC-SPR's contribution is methodological (interpretable sparse representatives, geometry control, stability diagnostics), not predictive dominance.

\item \textbf{LUAD is subsampled.}
The current LUAD evaluation uses 60 of approximately 500 available TCGA-LUAD samples.
Backend rankings may shift at full scale.

\item \textbf{PHATE-like backend is a lightweight surrogate.}
Our PHATE-like implementation computes potential distances without full multidimensional scaling; results may differ with the complete PHATE pipeline.

\item \textbf{Single-cell evidence is resource-bounded.}
The GSE165087 evaluation produces a ceiling effect (all methods tie at F1 $= 0.925$) due to the bounded cell count (320 cells).
Larger evaluations require memory-optimized implementations.

\item \textbf{Arabidopsis is a single cross-domain test.}
The plant evaluation demonstrates generalization to one non-cancer context; broader cross-domain evidence remains future work.

\item \textbf{Enrichment is exploratory.}
Pathway enrichment results are not corrected for multiple comparisons across datasets and should be interpreted as hypothesis-generating.

\item \textbf{Computational cost.}
GUDHI-based persistence computation scales quadratically with sample count, and Ricci curvature adds an iterative $O(n^2 k T)$ component.
The current CPU-only implementation limits applicability to datasets with $n < 5{,}000$ samples.
\end{enumerate}

% ============================================================
\section{Related Work}
\label{sec:related}
% ============================================================

\paragraph{Persistent homology in biology.}
TDA has been applied to gene expression \cite{rizvi2017}, protein structure \cite{xia2014}, tumor ecosystems \cite{vipond2021}, and single-cell trajectories \cite{moon2019phate}.
Most applications use standard Euclidean persistent homology without representative cycle analysis.

\paragraph{Harmonic persistent homology.}
The harmonic PH framework \cite{harmonicPH} introduced the theoretical bridge from persistent classes to sample-level representatives, and Damiani et al.\ \cite{damiani2024} demonstrated its application to omics data with breast cancer clustering.
CC-SPR builds directly on this work, preserving its scaffold while modifying the representative objective and generalizing the metric.

\paragraph{Ricci curvature on graphs.}
Ollivier--Ricci curvature \cite{ollivier2007, ollivier2009} and its discrete flow variant \cite{ni2019community} have been used for community detection and graph classification.
The connection to persistent homology through filtration perturbation appears novel.

\paragraph{Robust topology.}
DTM-based persistent homology \cite{chazal2011dtm, chazal2017robust} provides noise-robust topological inference.
Its integration into a representative-cycle framework with feature projection appears unexplored.

\paragraph{Sparse cycles.}
Sparse representations in TDA have been studied through $\ell_1$ optimal cycles \cite{chen2011hardness, dey2011optimal}, but typically without the biological feature-projection layer that CC-SPR provides.

\paragraph{Arabidopsis root development.}
Denyer et al.\ \cite{denyer2019} characterized spatiotemporal developmental trajectories in the Arabidopsis root using scRNA-seq, identifying cell-type-specific marker programs that we use as biological ground truth.

% ============================================================
\section{Conclusion}
\label{sec:conclusion}
% ============================================================

CC-SPR extends harmonic persistent homology with sparse elastic representatives and a backend-aware geometry framework.
Evaluation across cancer and plant transcriptomics demonstrates three findings:
(1)~backend ranking is dataset-dependent, confirming that geometry is a genuine modeling parameter;
(2)~topology-derived gene modules recover known biology with backend-specific emphasis (DTM: xenobiotic metabolism at $p < 0.001$; zero shared genes across all backends);
(3)~the framework generalizes across biological domains, with DTM outperforming Euclidean on both human cancer and Arabidopsis root development.

While standard baselines win on classification accuracy, CC-SPR provides complementary value through geometry-sensitive, stability-aware feature attribution.
The practical recommendation is to use CC-SPR not as a replacement for standard classifiers but as a complementary analysis tool that provides topology-grounded feature selection with explicit geometry control.

\paragraph{Data and Software Availability.}
CC-SPR is implemented as a Python package requiring Python $\geq$ 3.10 with dependencies including GUDHI, CVXPY, GraphRicciCurvature, and scanpy.
Source code, experiment scripts, and documentation are available at [GitHub URL].
All datasets are publicly available: TCGA-LUAD from the UCSC Xena Browser, GSE161711 and GSE165087 from NCBI GEO, and the Arabidopsis root atlas (GSE123013) from NCBI GEO.
Download scripts and reproducibility materials are provided in the repository.

\bibliographystyle{plain}
\begin{thebibliography}{25}

\bibitem{harmonicPH}
L.~Gurnari, A.~Monod, and S.~Oudot.
Harmonic persistent homology.
\emph{arXiv preprint arXiv:2405.16890}, 2024.

\bibitem{damiani2024}
D.~Damiani, O.~Vipond, and H.~A.~Harrington.
Probing omics data via harmonic persistent homology.
\emph{bioRxiv}, 2024.
\texttt{doi:10.1101/2024.09.15.613148}.

\bibitem{edelsbrunner2010}
H.~Edelsbrunner and J.~Harer.
\emph{Computational Topology: An Introduction}.
AMS, 2010.

\bibitem{carlsson2009}
G.~Carlsson.
Topology and data.
\emph{Bulletin of the AMS}, 46(2):255--308, 2009.

\bibitem{rizvi2017}
A.~H.~Rizvi, P.~G.~Camara, E.~K.~Kandber, et al.
Single-cell topological RNA-seq analysis reveals insights into cellular differentiation and development.
\emph{Nature Biotechnology}, 35(6):551--560, 2017.

\bibitem{xia2014}
K.~Xia and G.-W.~Wei.
Persistent homology analysis of protein structure, flexibility, and folding.
\emph{Int.\ J.\ Numer.\ Meth.\ Biomed.\ Eng.}, 30(8):814--844, 2014.

\bibitem{vipond2021}
O.~Vipond, J.~A.~Bull, P.~S.~Sherland, et al.
Multiparameter persistent homology landscapes identify immune cell spatial patterns in tumors.
\emph{PNAS}, 118(41), 2021.

\bibitem{moon2019phate}
K.~R.~Moon, D.~van Dijk, Z.~Wang, et al.
Visualizing structure and transitions in high-dimensional biological data.
\emph{Nature Biotechnology}, 37(12):1482--1492, 2019.

\bibitem{zomorodian2005}
A.~Zomorodian and G.~Carlsson.
Computing persistent homology.
\emph{Discrete \& Computational Geometry}, 33(2):249--274, 2005.

\bibitem{zou2005elasticnet}
H.~Zou and T.~Hastie.
Regularization and variable selection via the elastic net.
\emph{JRSS Series B}, 67(2):301--320, 2005.

\bibitem{ollivier2007}
Y.~Ollivier.
Ricci curvature of metric spaces.
\emph{C.\ R.\ Math.}, 345(11):643--646, 2007.

\bibitem{ollivier2009}
Y.~Ollivier.
Ricci curvature of Markov chains on metric spaces.
\emph{J.\ Funct.\ Anal.}, 256(3):810--864, 2009.

\bibitem{ni2019community}
C.-C.~Ni, Y.-Y.~Lin, F.~Luo, and J.~Gao.
Community detection on networks with Ricci flow.
\emph{Scientific Reports}, 9:9984, 2019.

\bibitem{coifman2006diffusion}
R.~R.~Coifman and S.~Lafon.
Diffusion maps.
\emph{Appl.\ Comput.\ Harm.\ Anal.}, 21(1):5--30, 2006.

\bibitem{chazal2011dtm}
F.~Chazal, D.~Cohen-Steiner, and Q.~M\'erigot.
Geometric inference for probability measures.
\emph{Found.\ Comput.\ Math.}, 11(6):733--751, 2011.

\bibitem{chazal2017robust}
F.~Chazal, B.~Fasy, F.~Lecci, B.~Michel, A.~Rinaldo, and L.~Wasserman.
Robust topological inference: distance to a measure and kernel distance.
\emph{JMLR}, 18(159):1--40, 2017.

\bibitem{gudhi2014}
The GUDHI Project.
GUDHI User and Reference Manual, 2014--2024.

\bibitem{chen2011hardness}
C.~Chen and D.~Freedman.
Hardness results for homology localization.
\emph{Discrete \& Comput.\ Geom.}, 45(3):425--448, 2011.

\bibitem{dey2011optimal}
T.~K.~Dey, J.~Sun, and Y.~Wang.
Approximating loops in a shortest homology basis from point data.
\emph{SoCG}, 2010.

\bibitem{denyer2019}
T.~Denyer, X.~Ma, S.~Klesen, et al.
Spatiotemporal developmental trajectories in the Arabidopsis root revealed using high-throughput single-cell RNA sequencing.
\emph{Developmental Cell}, 48(6):840--852, 2019.

\bibitem{gseapy}
Z.~Fang, X.~Liu, and G.~Peltz.
GSEApy: a comprehensive package for performing Gene Set Enrichment Analysis in Python.
\emph{Bioinformatics}, 39(1):btac757, 2023.

\end{thebibliography}

\end{document}
